\documentclass[11pt,a4paper]{article}
\usepackage{gtu-el-yazisi}
\begin{document}
\notbilgi{İşaretler ve Sistemler}{Ders 8 -- Sürekli zamanlı Fourier dönüşümü}{Isaret-ve-Sistemler-ders-8.pdf}
% Not: kaynak PDF, el yazısı notların ekran görüntülerinden oluşuyor; renkli kalemle sonradan eklenmiş işaretlerin (altı çizili yerler, oklar) bir kısmı vurgu olarak gösterildi.

\section{Sürekli Zamanlı Fourier Transformu (Dönüşümü)}
\sayfa{1}
(Continuous-Time Fourier Transform)

Periyodik olmayan (aperiyodik) işaretler için Fourier dönüşümü tanımlıdır.
\[ X(f)=\Fou{x(t)}=\int_{-\infty}^{\infty} x(t)\,e^{-j2\pi ft}\,dt \]
\[ x(t)=\iFou{X(f)}=\int_{-\infty}^{\infty} X(f)\,e^{j2\pi ft}\,df \]
\[ x(t)\ \underset{\mathcal{F}^{-1}}{\overset{\mathcal{F}}{\rightleftarrows}}\ X(f) \]
$\Rightarrow$ $\omega$ notasyonu:
\[ X(\omega)=\int_{-\infty}^{\infty} x(t)\,e^{-j\omega t}\,dt \]
\[ x(t)=\frac{1}{2\pi}\int_{-\infty}^{\infty} X(\omega)\,e^{j\omega t}\,d\omega, \qquad \omega=2\pi f=2\pi/T \]

$\Rightarrow$ Periyodik olmayan bir işaretin belli bir $T_0$ aralığında Fourier serisi açılımında $T_0$ sonsuza götürüldüğünde ($T_0\to\infty$) Fourier transformu elde edilir.
\begin{center}
\begin{tikzpicture}[x=0.6cm,y=0.8cm,murekkep]
  \draw[->] (-3.6,0) -- (3.4,0) node[right] {$t$};
  \draw[->] (0,-0.2) -- (0,1.6) node[right] {$x(t)$};
  \draw[thick] (-1,0) -- (-1,1) -- (1,1) -- (1,0);
  \node[font=\scriptsize,below] at (-1,0) {$-T_1$};
  \node[font=\scriptsize,below] at (1,0) {$T_1$};
  \draw[yesilkalem,thick] (-3,-0.2) -- (-3,0.2);
  \draw[yesilkalem,thick] (2.3,-0.2) -- (2.3,0.2);
  \node[yesilkalem,above] at (2.6,0.2) {$T_0$};
  \draw[yesilkalem,thick,->] (2.4,-0.5) -- (4.6,-0.5);
\end{tikzpicture}

\begin{tikzpicture}[x=0.6cm,y=0.8cm,murekkep]
  \draw[->] (-8.5,0) -- (8.5,0) node[right] {$t$};
  \draw[->] (0,-0.2) -- (0,1.6) node[right] {$x_{T_0}(t)$};
  \foreach \c in {-6,0,6} {\draw[thick] (\c-1,0) -- (\c-1,1) -- (\c+1,1) -- (\c+1,0);}
  \node at (-8.2,0.4) {$\cdots$}; \node at (8.1,0.4) {$\cdots$};
  \draw (-6,0.08) -- (-6,-0.08); \draw (6,0.08) -- (6,-0.08);
  \node[font=\scriptsize,below] at (-6,0) {$-T_0$};
  \node[font=\scriptsize,below] at (-1,0) {$-T_1$};
  \node[font=\scriptsize,below] at (1,0) {$T_1$};
  \node[font=\scriptsize,below] at (6,0) {$T_0$};
  \draw[yesilkalem,thick] (6,-0.1) .. controls (7,-0.35) and (7.6,-0.05) .. (8.6,-0.2);
\end{tikzpicture}
\end{center}
\vurgu[yesil]{$T_0\to\infty$} iken \vurgu[yesil]{$x_{T_0}(t)=x(t)$} dır (herhangi bir sonlu $t$ değeri için)
\[ {\color{yesilkalem}\longrightarrow}\quad \lim_{T_0\to\infty} x_{T_0}(t)=x(t) \]
\[ x_{T_0}(t)=\sum_{k=-\infty}^{\infty} x_k\,e^{jk\omega_0 t}, \qquad \omega_0=2\pi f_0=\frac{2\pi}{T_0} \]
\[ x_k=\frac{1}{T_0}\int_{-T_0/2}^{T_0/2} x_{T_0}(t)\,e^{-jk\omega_0 t}\,dt
\underset{\substack{\uparrow\\ x(t)=\left\{\begin{smallmatrix} x_{T_0}(t), & |t|<T_0/2\\ 0, & |t|>T_0/2\end{smallmatrix}\right.}}{=}
\frac{1}{T_0}\int_{-T_0/2}^{T_0/2} x(t)\,e^{-jk\omega_0 t}\,dt
\underset{\substack{\uparrow\\ T_0\to\infty}}{=}
\frac{1}{T_0}\int_{-\infty}^{\infty} x(t)\,e^{-jk\omega_0 t}\,dt \]

\sayfa{2}
\noindent
\begin{minipage}[t]{0.47\linewidth}
\[ {\color{yesilkalem}\longrightarrow}\ X(f)=\int_{-\infty}^{\infty} x(t)\,e^{-j2\pi ft}\,dt \]
\[ {\color{yesilkalem}\longrightarrow}\ x_k=\frac{1}{T_0}X(kf_0)\ {\color{yesilkalem}\longleftarrow} \]
\[ x_{T_0}(t)=\sum_{k=-\infty}^{\infty}\frac{1}{T_0}X(kf_0)\,e^{j2\pi kf_0t} \]
\[ {\color{yesilkalem}\longrightarrow}\ \Delta f=\frac{1}{T_0} \]
\begin{center}
\vurgu[yesil]{$\displaystyle x_{T_0}(t)=\sum_{k=-\infty}^{\infty} X(k\Delta f)\,e^{j2\pi k\Delta ft}\,\Delta f$}
\end{center}
\begin{align*}
x(t) &= \lim_{T_0\to\infty} x_{T_0}(t)\\
&= \underbrace{\lim_{\Delta f\to0}\sum_{k=-\infty}^{\infty} X(k\Delta f)\,e^{j2\pi k\Delta ft}\,\Delta f}_{X(f)e^{j2\pi ft}\text{'nin altında kalan alan}}
\ {\color{yesilkalem}\big\downarrow}
\end{align*}
\[ x(t)=\int_{-\infty}^{\infty} X(f)\,e^{j2\pi ft}\,df \]
\end{minipage}\hfill
\begin{minipage}[t]{0.06\linewidth}
\vspace{18mm}\centering
{\color{kirmizikalem}\fbox{\scriptsize$\begin{array}{@{}c@{}}df=1/T_0\\=f_0\end{array}$}}
\end{minipage}\hfill
\begin{minipage}[t]{0.44\linewidth}
\[ X(\omega)=\int_{-\infty}^{\infty} x(t)\,e^{-j\omega t} \]
\[ x_k=\frac{1}{T_0}X(k\omega_0) \]
\[ x_{T_0}(t)=\sum_{k=-\infty}^{\infty}\frac{1}{T_0}X(k\omega_0)\,e^{jk\omega_0t} \]
\[ \Delta\omega=\frac{2\pi}{T_0} \]
\[ x_{T_0}(t)=\frac{1}{2\pi}\sum_{k=-\infty}^{\infty} X(k\Delta\omega)\,e^{jk\Delta\omega t}\,\Delta\omega \]
\[ x(t)=\underbrace{\lim_{\Delta\omega\to0}\frac{1}{2\pi}\sum_{k=-\infty}^{\infty} X(k\Delta\omega)\,e^{jk\Delta\omega t}\,\Delta\omega}_{X(\omega)e^{j\omega t}} \]
\[ x(t)=\int_{-\infty}^{\infty} X(\omega)\,e^{j\omega t}\,d\omega\quad\text{\footnotemark} \]
\end{minipage}
\footnotetext{Claude notu: muhtemelen önünde $\frac{1}{2\pi}$ olmalı: $x(t)=\frac{1}{2\pi}\int_{-\infty}^{\infty} X(\omega)e^{j\omega t}d\omega$.}

\medskip
$\Rightarrow$ Fourier dönüşümünün mevcut olma koşulları (Dirichlet condition)
\begin{enumerate}[label=\arabic*)]
\item Aperyodik $x(t)$ işareti \vurgu[pembe]{mutlak integrallenebilir} olmalıdır, \ $\displaystyle\int_{-\infty}^{\infty} |x(t)|\,dt<\infty$
\item \vurgu[pembe]{Herhangi bir sonlu $t$ aralığında},
  \begin{itemize}[label=--]
  \item sonlu sayıda \vurgu[pembe]{maksimum ve minimum} değere sahip olmalıdır.
  \item sonlu sayıda \vurgu[pembe]{süreksizliğe} ve bu noktalarda sonlu \vurgu[pembe]{değere sahip} olmalıdır.
  \end{itemize}
\end{enumerate}

$\Rightarrow$ Frekans spektrumları:
\[ {\color{kirmizikalem}\nearrow}\ X(f)=|X(f)|\,e^{j\theta_x(f)} \;\longrightarrow\;
\begin{aligned}[t] &|X(f)|:\ \text{genlik spektrumu}\\ &\hookrightarrow \theta_x(f):\ \text{faz spektrumu}\end{aligned} \]

\sayfa{3}
\begin{ornek}
Ex: $x(t)=e^{-at}u(t)$, \ $a>0$ ise \vurgu[turuncu]{$X(f)=\,?$} \ Frekans spektrumlarını çiziniz.
\end{ornek}
\[ X(f)=\int_{-\infty}^{\infty} x(t)\,e^{-j2\pi ft}\,dt=\int_{0}^{\infty} e^{-at}\cdot e^{-j2\pi ft}\,dt=\int_{0}^{\infty} e^{-(a+j2\pi f)t}\,dt=\left.\frac{e^{-(a+j2\pi f)t}}{-(a+j2\pi f)}\right|_0^{\infty}=\frac{1}{a+j2\pi f}=\frac{1}{a+j\omega} \]
\[ |X(f)|=\frac{1}{\sqrt{a^2+(2\pi f)^2}} \]
\[ \theta_x(f)=-\tan^{-1}(2\pi f/a) \]
\begin{center}
\begin{tikzpicture}
\begin{axis}[width=6.4cm,height=4.2cm,axis lines=middle,xmin=-0.5,xmax=0.55,ymin=0,ymax=1.3,
    xtick=\empty,ytick={1},yticklabels={$1/a$},xlabel={$f$},ylabel={$|X(f)|$},
    every axis x label/.style={at={(current axis.right of origin)},anchor=west},
    every axis y label/.style={at={(current axis.above origin)},anchor=south}]
  \addplot[domain=-0.5:0.5,samples=80,smooth,thick,murekkep] {1/sqrt(1+(2*pi*x)^2)};
\end{axis}
\end{tikzpicture}\hspace{1cm}
\begin{tikzpicture}
\begin{axis}[width=6.4cm,height=4.6cm,axis lines=middle,xmin=-1,xmax=1.1,ymin=-1.3,ymax=1.4,
    xtick=\empty,ytick={-1,1},yticklabels={$-\pi/2$,$\pi/2$},ylabel={$\theta_x(f)$},
    every axis y label/.style={at={(current axis.above origin)},anchor=south}]
  \addplot[domain=-1:1,samples=80,smooth,thick,murekkep] {-atan(2*pi*x)/90};
  \addplot[dashed,kursunkalem,domain=-1:1] {1};
  \addplot[dashed,kursunkalem,domain=-1:1] {-1};
\end{axis}
\end{tikzpicture}
\end{center}

\begin{ornek}
Ex: $x(t)=e^{-a|t|}$, \ $a>0$ ise $X(f)=\,?$
\end{ornek}
\begin{align*}
X(f) &= \int_{-\infty}^{\infty} x(t)\,e^{-j2\pi ft}\,dt
\underset{\substack{\uparrow\\ x(t)=\left\{\begin{smallmatrix} e^{at}, & t<0\\ e^{-at}, & t>0\end{smallmatrix}\right.}}{=}
\int_{-\infty}^{0} e^{at}e^{-j2\pi ft}\,dt+\int_{0}^{\infty} e^{-at}e^{-j2\pi ft}\,dt\\
&= \left.\frac{e^{(a-j2\pi f)t}}{(a-j2\pi f)}\right|_{-\infty}^{0}+\left.e^{-(a+j2\pi f)t}\right|_{0}^{\infty}
=\frac{1}{a-j2\pi f}+\frac{1}{a+j2\pi f}=\frac{2a}{a^2+(2\pi f)^2}\quad\text{\footnotemark}
\end{align*}
\footnotetext{Claude notu: ikinci terimde $-(a+j2\pi f)$ paydası unutulmuş; sonuç doğru.}

\noindent
\begin{minipage}[c]{0.4\linewidth}
\[ |X(f)|=\frac{2a}{a^2+(2\pi f)^2} \]
\[ \theta_x(f)=0 \]
\end{minipage}\hfill
\begin{minipage}[c]{0.56\linewidth}
\centering
\begin{tikzpicture}
\begin{axis}[width=7cm,height=4.2cm,xmin=-0.5,xmax=0.55,ymin=0,ymax=2.6,
    xtick=\empty,ytick={2},yticklabels={$2/a$},xlabel={$f$},ylabel={$X(f)$}]
  \addplot[domain=-0.5:0.5,samples=80,smooth,thick,murekkep] {2/(1+(2*pi*x)^2)};
\end{axis}
\end{tikzpicture}
\end{minipage}



\sayfa{4}
\begin{ornek}
Ex: \vurgu[turuncu]{$x(t)$}$=\begin{cases} 1, & |t|<\tau/2\\ 0, & \text{dışında}\end{cases}$ \quad \vurgu[turuncu]{dikdörtgen darbe işaretinin} FT bulunuz.
\end{ornek}
\begin{center}
\begin{tikzpicture}[x=1cm,y=0.8cm,murekkep]
  \draw[->] (-1.6,0) -- (1.8,0) node[right] {$t$};
  \draw[->] (0,-0.1) -- (0,1.6) node[right] {$x(t)=\Pi\!\left(\frac{t}{\tau}\right)$};
  \draw[thick] (-0.8,0) -- (-0.8,1) -- (0.8,1) -- (0.8,0);
  \node[font=\scriptsize,above left] at (0,1) {$1$};
  \node[font=\scriptsize,below] at (-0.8,0) {$-\tau/2$};
  \node[font=\scriptsize,below] at (0.8,0) {$\tau/2$};
\end{tikzpicture}
\end{center}
\begin{align*}
X(f) &= \int_{-\infty}^{\infty} x(t)\,e^{-j2\pi ft}\,dt=\int_{-\tau/2}^{\tau/2} e^{-j2\pi ft}\,dt=\left.\frac{e^{-j2\pi ft}}{-j2\pi f}\right|_{-\tau/2}^{\tau/2}\\
&= \frac{-1}{j2\pi f}\left[e^{-j2\pi f\tau/2}-e^{+j2\pi f\tau/2}\right]=\frac{1}{\pi f}\sin(\pi f\tau)
\end{align*}

Reminder: \quad \vurgu[turuncu]{$\operatorname{sinc}(t)$}$=\begin{cases} \dfrac{\sin(\pi t)}{\pi t}, & t\neq0\\[2mm] 1, & t=0\end{cases}$
\begin{center}
\begin{tikzpicture}
\begin{axis}[width=8cm,height=4.4cm,xmin=-3.6,xmax=3.9,ymin=-0.35,ymax=1.25,
    xtick={-3,-2,-1,1,2,3},ytick={1},xlabel={$t$},ylabel={$\operatorname{sinc}(t)$}]
  \addplot[domain=-3.5:3.5,samples=200,smooth,thick,murekkep] {sin(deg(pi*x))/(pi*x)};
  \addplot[only marks,mark=*,mark size=1.5pt,turuncukalem] coordinates {(0,0) (1,0) (2,0) (3,0)};
  \node at (axis cs:-3.4,0.12) {$\cdots$}; \node at (axis cs:3.6,-0.12) {$\cdots$};
\end{axis}
\end{tikzpicture}
\end{center}
\[ X(f)=\frac{1}{\pi f}\sin(\pi f\tau)=\tau\,\textcolor{turuncukalem}{\boxed{\textcolor{murekkep}{\frac{\sin(\pi f\tau)}{\pi f\tau}}}}=\tau\operatorname{sinc}(f\tau) \]

\noindent
\begin{minipage}[c]{0.58\linewidth}
\centering
\begin{tikzpicture}
\begin{axis}[width=9.4cm,height=4.6cm,xmin=-4.4,xmax=4.7,ymin=-0.35,ymax=1.25,
    xtick={-3,-2,-1,1,2,3,4},xticklabels={$-3/\tau$,$-2/\tau$,$-1/\tau$,$1/\tau$,$2/\tau$,$3/\tau$,$4/\tau$},
    ytick=\empty,xlabel={$f$},ylabel={$X(f)$},tick label style={font=\scriptsize}]
  \addplot[domain=-4.2:4.3,samples=200,smooth,thick,murekkep] {sin(deg(pi*x))/(pi*x)};
  \addplot[only marks,mark=*,mark size=1.5pt,turuncukalem] coordinates {(-1,0) (1,0) (2,0) (3,0) (4,0) (0,1)};
  \node at (axis cs:-4.1,0.15) {$\cdots$}; \node at (axis cs:4.4,0.2) {$\cdots$};
\end{axis}
\end{tikzpicture}
\end{minipage}\hfill
\begin{minipage}[c]{0.40\linewidth}
\centering
$x(t)={\color{turuncukalem}A}\,\Pi\!\left(\dfrac{t}{\tau}\right) \xleftrightarrow{\ \mathcal{F}\ } X(f)={\color{turuncukalem}A}\,\tau\operatorname{sinc}(f\tau)$

\smallskip
\begin{tikzpicture}[x=1cm,y=0.7cm,murekkep]
  \draw[->] (-1.4,0) -- (1.6,0) node[right] {$t$};
  \draw[->] (0,-0.1) -- (0,1.6) node[right] {$x(t)$};
  \draw[thick] (-0.7,0) -- (-0.7,1) -- (0.7,1) -- (0.7,0);
  \node[font=\scriptsize,above left] at (0,1) {$1$};
  \node[font=\scriptsize,below] at (-0.7,0) {$-\tau/2$};
  \node[font=\scriptsize,below] at (0.7,0) {$\tau/2$};
\end{tikzpicture}
\end{minipage}

\sayfa{5}
\begin{ornek}
Ex: \vurgu[turuncu]{$x(t)=\delta(t)$} ise \vurgu[turuncu]{$X(f)=\,?$}
\end{ornek}
\[ X(f)=\int_{-\infty}^{\infty}\delta(t)\,e^{-j2\pi ft}\,dt=\int_{-\infty}^{\infty}\delta(t)\,\overset{\color{turuncukalem}1}{e^{-j2\pi f\cdot 0}}\,dt=\int_{-\infty}^{\infty}\delta(t)\,dt=1 \]
\[ x(t)=\delta(t) \xleftrightarrow{\ \mathcal{F}\ } X(f)=1 \]
\begin{center}
\begin{tikzpicture}[x=1cm,y=0.8cm,murekkep]
  \draw[->] (-1.5,0) -- (1.8,0) node[right] {$t$};
  \draw[->] (0,-0.1) -- (0,1.7) node[right] {$x(t)$};
  \draw[->,very thick,turuncukalem] (0,0) -- (0,1);
  \node[font=\scriptsize,left] at (0,0.9) {$1$};
\end{tikzpicture}
\hspace{2cm}
\begin{tikzpicture}[x=1cm,y=0.8cm,murekkep]
  \draw[->] (-1.8,0) -- (1.8,0) node[right] {$f$};
  \draw[->] (0,-0.1) -- (0,1.7) node[right] {$X(f)$};
  \draw[thick] (-1.6,1) -- (1.6,1);
  \node[font=\scriptsize,above right] at (0,1) {$1$};
\end{tikzpicture}
\end{center}

\begin{ornek}
Ex: $X(f)=\delta(f-f_0)$ ise $x(t)=\,?$
\end{ornek}
\[ x(t)=\int_{-\infty}^{\infty} X(f)\,e^{j2\pi ft}\,df=\int_{-\infty}^{\infty}\delta(f-f_0)\,e^{j2\pi ft}\,df=\int_{-\infty}^{\infty}\delta(f-f_0)\,e^{j2\pi f_0t}\,df=e^{j2\pi f_0t}\underbrace{\int_{-\infty}^{\infty}\delta(f-f_0)\,df}_{=1}=e^{j2\pi f_0t} \]
\begin{ornek}
Ex: \vurgu[turuncu]{$X(f)=\delta(f)$} ise \vurgu[turuncu]{$x(t)=\,?$}
\end{ornek}
\[ x(t)=\int_{-\infty}^{\infty} X(f)\,e^{j2\pi ft}\,df=\int_{-\infty}^{\infty}\delta(f)\,e^{j2\pi ft}\,df=\int_{-\infty}^{\infty}\delta(f)\,e^{j2\pi 0t}\,df=\int_{-\infty}^{\infty}\delta(f)\,df=1 \]
\begin{align*}
\delta(t) &\xleftrightarrow{\ \mathcal{F}\ } 1\\
1 &\xleftrightarrow{\ \mathcal{F}\ } \delta(f)\\
e^{j2\pi f_0t} &\xleftrightarrow{\ \mathcal{F}\ } \delta(f-f_0)\\
e^{-j2\pi f_0t} &\xleftrightarrow{\ \mathcal{F}\ } \delta(f+f_0)
\end{align*}

\sayfa{6}
\section{Periyodik işaretlerin Fourier Dönüşümü}
\[ x(t)=x(t+T),\ \forall t \qquad T=\frac{1}{f_0}=\frac{2\pi}{\omega_0} \]
\[ x(t)=\sum_{k=-\infty}^{\infty} x_k\,e^{jk\omega_0t}\underset{\substack{\uparrow\\ \omega_0=2\pi f_0}}{=}\sum_{k=-\infty}^{\infty} x_k\,e^{j2\pi kf_0t}=\cdots+x_{-1}e^{-j2\pi f_0t}+x_0+x_1e^{j2\pi f_0t}+x_2e^{j2\pi 2f_0t}+\cdots \]
\begin{align*}
x_1e^{+j2\pi f_0t} &\xleftrightarrow{\ \mathcal{F}\ } x_1\,\delta(f-f_0)\\
x_0 &\xleftrightarrow{\ \mathcal{F}\ } x_0\,\delta(f)\\
x_2e^{j2\pi 2f_0t} &\xleftrightarrow{\ \mathcal{F}\ } x_2\,\delta(f-2f_0)
\end{align*}
\[ x(t)=\sum_{k=-\infty}^{\infty} x_k\,e^{j2\pi kf_0t} \xleftrightarrow{\ \mathcal{F}\ } X(f)=\sum_{k=-\infty}^{\infty} x_k\,\delta(f-kf_0) \]
\[ x_k=\frac{1}{T_0}\int_{T_0} x(t)\,e^{-jk\omega_0t}\,dt \]

\begin{ornek}
Ex: $x(t)=\cos 2\pi f_0t$ ise $X(f)=\,?$
\end{ornek}
\[ x(t)=\frac{1}{2}\left[e^{j2\pi f_0t}+e^{-j2\pi f_0t}\right] \;\longrightarrow\; x_1=x_{-1}=1/2 \]
\[ X(f)=\frac{1}{2}\delta(f-f_0)+\frac{1}{2}\delta(f+f_0) \]
\begin{center}
\begin{tikzpicture}[x=1cm,y=1cm,murekkep]
  \draw[->] (-2.2,0) -- (2.4,0) node[right] {$f$};
  \draw[->] (0,-0.1) -- (0,1.5) node[right] {$X(f)$};
  \draw[->,very thick] (-1.2,0) -- (-1.2,0.8);
  \draw[->,very thick] (1.2,0) -- (1.2,0.8);
  \draw (-0.06,0.8) -- (0.06,0.8) node[right,font=\scriptsize] {$1/2$};
  \node[font=\scriptsize,below] at (-1.2,0) {$-f_0$};
  \node[font=\scriptsize,below] at (1.2,0) {$f_0$};
\end{tikzpicture}
\end{center}

\begin{ornek}
Ex: $x(t)=\sin 2\pi f_0t$
\end{ornek}
\[ \begin{aligned}
x(t) &= \frac{1}{2j}\left[e^{j2\pi f_0t}-e^{-j2\pi f_0t}\right]\\
X(f) &= \frac{1}{2j}\left[\delta(f-f_0)-\delta(f+f_0)\right]
\end{aligned}
\quad
\begin{aligned}
&\longrightarrow x_1=\frac{1}{2j}=\frac{-1}{2}j\\
&\phantom{\longrightarrow\ } x_{-1}=-\frac{1}{2j}=\frac{1}{2}j
\end{aligned}
\ \left.\vphantom{\begin{aligned}a\\[4mm]a\end{aligned}}\right\}\quad
\begin{aligned}
&|x_1|=|x_{-1}|=1/2\\
&\theta_{x_1}=-\pi/2 \quad \theta_{x_{-1}}=\pi/2
\end{aligned} \]
\begin{center}
\begin{tikzpicture}[x=1cm,y=1cm,murekkep,baseline=(current bounding box.center)]
  \draw[->] (-2,0) -- (2.2,0) node[right] {$f$};
  \draw[->] (0,-1.2) -- (0,1.4) node[right] {$X(f)$};
  \draw[->,very thick] (1.1,0) -- (1.1,0.8) node[right,font=\scriptsize] {$1/2j$};
  \draw[->,very thick] (-1.1,0) -- (-1.1,-0.8) node[right,font=\scriptsize] {$-1/2j$};
  \node[font=\scriptsize,above] at (-1.1,0) {$-f_0$};
  \node[font=\scriptsize,below] at (1.1,0) {$f_0$};
\end{tikzpicture}
\hspace{1.2cm}
\begin{tikzpicture}[x=1cm,y=1cm,murekkep,baseline=(current bounding box.center)]
  \begin{scope}[yshift=1.6cm]
  \draw[->] (-2,0) -- (2.2,0) node[right] {$f$};
  \draw[->] (0,-0.1) -- (0,1.3) node[right] {$|X(f)|$};
  \draw[->,very thick] (-1.1,0) -- (-1.1,0.8);
  \draw[->,very thick] (1.1,0) -- (1.1,0.8);
  \draw (-0.06,0.8) -- (0.06,0.8) node[right,font=\scriptsize] {$1/2$};
  \node[font=\scriptsize,below] at (-1.1,0) {$-f_0$};
  \node[font=\scriptsize,below] at (1.1,0) {$f_0$};
  \end{scope}
  \begin{scope}[yshift=-0.6cm]
  \draw[->] (-2,0) -- (2.2,0) node[right] {$f$};
  \draw[->] (0,-1) -- (0,1.2) node[right] {$\theta_x(f)$};
  \draw[very thick] (-1.1,0) -- (-1.1,0.8) node[left,font=\scriptsize] {$\pi/2$};
  \draw[very thick] (1.1,0) -- (1.1,-0.8); \fill (1.1,-0.8) circle (1.5pt) node[right,font=\scriptsize] {$-\pi/2$};
  \node[font=\scriptsize,below left] at (-1.1,0) {$-f_0$};
  \node[font=\scriptsize,above] at (1.1,0) {$f_0$};
  \end{scope}
\end{tikzpicture}
\end{center}


\sayfa{7}
\begin{ornek}
Ex: $x(t)=\displaystyle\sum_{k=-\infty}^{\infty}\delta(t-kT_0)$ dürtü katarı işaretinin FT bulunuz.
\end{ornek}
\noindent
\begin{minipage}[c]{0.45\linewidth}
\centering
\begin{tikzpicture}[x=0.75cm,y=1cm,murekkep]
  \draw[->] (-3,0) -- (3.2,0) node[right] {$t$};
  \draw[->] (0,-0.1) -- (0,1.3) node[right] {$x(t)$};
  \foreach \k in {-2,-1,0,1,2} {\draw[->,very thick] (\k,0) -- (\k,0.8);}
  \node at (-2.7,0.5) {$\cdots$}; \node at (2.7,0.5) {$\cdots$};
  \node[font=\scriptsize,below] at (-2,0) {$-2T_0$};
  \node[font=\scriptsize,below] at (-1,0) {$-T_0$};
  \node[font=\scriptsize,below] at (0,0) {$0$};
  \node[font=\scriptsize,below] at (1,0) {$T_0$};
  \node[font=\scriptsize,below] at (2,0) {$2T_0$};
\end{tikzpicture}
\end{minipage}\hfill
\begin{minipage}[c]{0.53\linewidth}
Reminder: \ $x_k=\dfrac{1}{T_0}\displaystyle\int_{T_0} x(t)\,e^{-jk\omega_0t}=\dfrac{1}{T_0}$
\end{minipage}
\[ X(f)=\sum_{k=-\infty}^{\infty} x_k\,\delta(f-kf_0)=\frac{1}{T_0}\sum_{k=-\infty}^{\infty}\delta(f-kf_0) \]
\begin{center}
\begin{tikzpicture}[x=0.75cm,y=1cm,murekkep]
  \draw[->] (-3.4,0) -- (4,0) node[right] {$f$};
  \draw[->] (0,-0.1) -- (0,1.3) node[right] {$X(f)$};
  \foreach \k in {-2,-1,0,1,2,3} {\draw[->,very thick] (\k,0) -- (\k,0.8);}
  \node at (-3,0.5) {$\cdots$}; \node at (3.6,0.5) {$\cdots$};
  \node[font=\scriptsize,right] at (0,0.95) {$1/T_0$};
  \node[font=\scriptsize,below] at (-1,0) {$-f_0$};
  \node[font=\scriptsize,below] at (0,0) {$0$};
  \node[font=\scriptsize,below] at (1,0) {$f_0$};
  \node[font=\scriptsize,below] at (2,0) {$2f_0$};
\end{tikzpicture}
\end{center}

\sekilnotu[7]{Sayfanın sağ üst köşesinde turuncu kalemle eklenmiş kaba bir çizim: yatay bir eksen üzerinde aralıklı dikey çizgiler (dürtü katarına benzer).}

\section{Sürekli Zamanlı Fourier Transformunun Özellikleri}
\[ \begin{aligned} &x(t)\xleftrightarrow{\ \mathcal{F}\ } X(f)\\ &X(f)=\Fou{x(t)}\end{aligned} \qquad,\ x(t):\ \text{aperiyodik işaret} \]

1) Lineerlik: \quad $\alpha x_1(t)+\beta x_2(t) \xleftrightarrow{\ \mathcal{F}\ } \alpha X_1(f)+\beta X_2(f)$

2) Zamanda öteleme: \quad $x(t-t_0) \xleftrightarrow{\ \mathcal{F}\ } e^{-j2\pi ft_0}X(f)$
\[ \Fou{x(t-t_0)}=\int_{-\infty}^{\infty} x(t-t_0)\,e^{-j2\pi ft}\,dt\underset{\substack{\uparrow\\ t-t_0=\tau}}{=}\int_{-\infty}^{\infty} x(\tau)\,e^{-j2\pi f(\tau+t_0)}\,d\tau=e^{-j2\pi ft_0}\underbrace{\int_{-\infty}^{\infty} x(\tau)\,e^{-j2\pi f\tau}\,d\tau}_{X(f)} \]
Aperiyodik bir $x(t)$ işareti zamanda ötelendiğinde genlik spektrumu değişmez, faz spektrumu $-2\pi ft_0$ kadar ötelenir.

3) Frekansta öteleme: \quad $x(t)\,e^{j2\pi f_0t} \xleftrightarrow{\ \mathcal{F}\ } X(f-f_0)$

İspat:
\[ \Fou{x(t)e^{j2\pi f_0t}}=\int_{-\infty}^{\infty} x(t)\,e^{j2\pi f_0t}\,e^{-j2\pi ft}\,dt=\int_{-\infty}^{\infty} x(t)\,e^{-j2\pi(f-f_0)t}\,dt=X(f-f_0) \]
$x(t)$ işareti $e^{j2\pi f_0t}$ ile çarpıldığında frekans spektrumu $f_0$ kadar ötelenir.

4) Modülasyon: \quad $x(t)\cos 2\pi f_0t \xleftrightarrow{\ \mathcal{F}\ } \dfrac{1}{2}\left[X(f-f_0)+X(f+f_0)\right]$
\[ x(t)\sin 2\pi f_0t \xleftrightarrow{\ \mathcal{F}\ } \frac{1}{2j}\left[X(f-f_0)-X(f+f_0)\right]=\frac{1}{2}\left[X(f-f_0)\,e^{-j\pi/2}+X(f+f_0)\,e^{j\pi/2}\right] \]

\sayfa{8}
5) Zamanda ölçeklendirme: \quad $x(at) \xleftrightarrow{\ \mathcal{F}\ } \dfrac{1}{|a|}X(f/a)$, \ $a\neq0$
\quad $\longrightarrow$ \begin{tabular}[t]{@{}l@{}} $a>1$ işaret frekansta (spektrum) genişler.\\ $0<a<1$ işaret frekansta daralır.\end{tabular}

6) Duality: \quad $\begin{aligned}[t] x(t) &\xleftrightarrow{\ \mathcal{F}\ } X(f)\\ X(t) &\xleftrightarrow{\ \mathcal{F}\ } x(-f)\end{aligned}$

İspat:
\[ x(t)=\iFou{X(f)}=\int_{-\infty}^{\infty} X(f)\,e^{j2\pi ft}\,df \]
\[ x(-t)=\int_{-\infty}^{\infty} X(f)\,e^{-j2\pi ft}\,df\ \underset{\substack{\uparrow\\ t\to f}}{\Rightarrow}\ x(-f)=\textcolor{turuncukalem}{\boxed{\textcolor{murekkep}{\int_{-\infty}^{\infty} X(t)\,e^{-j2\pi ft}\,dt}}}\qquad {\color{turuncukalem}X(t)} \]

7) Türev alma: \quad $\dfrac{dx(t)}{dt} \xleftrightarrow{\ \mathcal{F}\ } j2\pi f\,X(f)$

\begin{ornek}
Ex: $x(t)=A\,\Pi(t/\tau) \xleftrightarrow{\ \mathcal{F}\ } X(f)=A\tau\operatorname{sinc}(f\tau)$ ise $x(t)=\operatorname{sinc}(t\tau)$ işaretinin $\text{FT}=\,?$
\end{ornek}
Reminder: \quad $\left.\begin{aligned} x(t)&\longleftrightarrow X(f)\\ X(f)&\longleftrightarrow x(-t)\end{aligned}\right\}$ Duality
\[ x(t)=A\tau\operatorname{sinc}(t\tau) \xleftrightarrow{\ \mathcal{F}\ } X(f)=A\,\Pi(-f/\tau)\underset{\substack{\uparrow\\ \text{Çift fonksiyon}}}{=}A\,\Pi(f/\tau) \]
\[ x(t)=\operatorname{sinc}(t\tau) \xleftrightarrow{\ \mathcal{F}\ } X(f)=\frac{1}{\tau}\,\Pi(f/\tau) \]

\sekilnotu[8]{Sağ tarafta turuncu kalemle çizilmiş kaba bir sinc eğrisi: orijinde tepe yapan, iki yana doğru sönümlenerek salınan dalga.}

\begin{ornek}
Ex: $x(t)=\dfrac{2}{1+t^2}$ ise $X(f)=\,?$
\end{ornek}
Reminder: \quad $e^{-a|t|} \longleftrightarrow \dfrac{2a}{a^2+(2\pi f)^2}$
\[ x(t)=e^{-|t|} \longleftrightarrow X(f)=\frac{2}{1+(2\pi f)^2} \]
\[ x(t)=\frac{2}{1+t^2} \longleftrightarrow X(f)=e^{-|-f|}=e^{-|f|}\quad\text{\footnotemark} \]
\footnotetext{Claude notu: $X(f)=\frac{2}{1+(2\pi f)^2}$ olduğundan dualite $\frac{2}{1+(2\pi t)^2}\leftrightarrow e^{-|f|}$ verir; $\frac{2}{1+t^2}$ için sonuç muhtemelen $2\pi e^{-2\pi|f|}$ olmalı.}


\sayfa{9}
\begin{ornek}
Ex: \vurgu[turuncu]{$x(t)=e^{-|t-t_0|}$} ise $X(f)=\,?$
\end{ornek}
Reminder: \quad $\textcolor{turuncukalem}{\boxed{\textcolor{murekkep}{x(t-t_0) \longleftrightarrow e^{-j2\pi ft_0}X(f)}}}$ \quad ve \quad $\textcolor{turuncukalem}{\boxed{\textcolor{murekkep}{e^{-|t|} \longleftrightarrow \dfrac{2a}{a^2+(2\pi f)^2}}}}$
\[ x(t)=e^{-|t-t_0|} \xleftrightarrow{\ \mathcal{F}\ } \textcolor{turuncukalem}{\boxed{\textcolor{murekkep}{X(f)=\frac{2a}{a^2+(2\pi f)^2}\,e^{-j2\pi ft_0}}}}
\;\longrightarrow\; \begin{aligned}[t] &|X(f)|=\frac{2a}{a^2+(2\pi f)^2}\\ &\hookrightarrow \theta_x(f)=-2\pi ft_0\end{aligned} \]
\begin{center}
\begin{tikzpicture}
\begin{axis}[width=6.6cm,height=4.2cm,xmin=-0.5,xmax=0.55,ymin=0,ymax=2.6,
    xtick=\empty,ytick={2},yticklabels={$2/a$},xlabel={$f$},ylabel={$|X(f)|$}]
  \addplot[domain=-0.5:0.5,samples=80,smooth,thick,murekkep] {2/(1+(2*pi*x)^2)};
\end{axis}
\end{tikzpicture}\hspace{1cm}
\begin{tikzpicture}[x=1cm,y=1cm,murekkep]
  \draw[->] (-1.6,0) -- (1.8,0) node[right] {$f$};
  \draw[->] (0,-1.5) -- (0,1.6) node[right] {$\theta_x(f)$};
  \draw[thick] (-1.3,1.2) -- (1.3,-1.2);
  \draw[turuncukalem,line width=2pt,opacity=0.5] (-1.3,1.2) -- (1.3,-1.2);
  \node[font=\scriptsize,right] at (1.3,-1.2) {$-2\pi ft_0$};
\end{tikzpicture}
\end{center}

\end{document}
